% section: 解法の地図

\section{解法の地図}

第2章では,漸化式をいくつかの型に分け,それぞれの型を解く方法を
見てきた.第3章では,個々の公式をもう一度並べるのではなく,
「式のどこを見て,どの操作を選ぶか」を一つの流れとして整理する.

漸化式を解くときの目標は,最初から一般項を当てることではない.
次の項を作る規則を,すでに知っている単純な規則へ翻訳することである.
その翻訳には,大きく次の道筋がある.

\begin{enumerate}
  \item \textbf{足し算の構造を見る.}
    差 $a_{n+1}-a_n$ が $n$ だけで表されるなら,差を足し戻す.
  \item \textbf{掛け算の構造を見る.}
    比 $a_{n+1}/a_n$ が分かるなら,積を作る.べきが現れるなら対数で
    掛け算を足し算へ移す.
  \item \textbf{定数を消す.}
    $a_{n+1}=pa_n+q$ のような式では,不動点を引いて等比型にする.
  \item \textbf{過去の項をまとめる.}
    $a_{n+2},a_{n+1},a_n$ が線形に現れるなら,特性方程式で基本解を
    探す.重解や複素数解も,同じ考えの続きである.
  \item \textbf{表現を変える.}
    逆数,一次分数変換,三角関数,対称式,母関数など,次の操作をしても
    同じ種類の表現が保たれる変数を探す.
  \item \textbf{一般項以外の情報を調べる.}
    閉じた式が得られないときは,不動点,単調性,有界性,極限,近似誤差を
    調べる.「解く」を一つの式を書くことに限定しない.
\end{enumerate}

この順序は,絶対的なアルゴリズムではない.一つの漸化式が複数の顔を
持つこともある.例えば,線形漸化式は特性方程式でも行列でも扱えるし,
三角関数型の漸化式はチェビシェフ多項式としても見られる.大切なのは,
式を見た瞬間に方法を決めつけず,どの表現なら規則が簡単になるかを
比較することである.

\begin{quote}
解法を選ぶときは,「この漸化式は何型か」と名前を当てるだけで終わらせない.
「何を単純にしたいのか」を先に決める.差を単純にしたいのか,比を単純に
したいのか,定数項を消したいのか,過去の項をまとめたいのか.
目標が決まれば,必要な変形はかなり自然に見えてくる.
\end{quote}

\subsection{解の形を三段階に分ける}

漸化式の「解」には,少なくとも三つの段階がある.

\begin{enumerate}
  \item \textbf{閉じた式:} $a_n$ を $n$ と初期値だけで直接表す.
  \item \textbf{構造を保った表現:} $a_n$ を積,和,行列,関数の反復などで
    表し,その表現から性質を読み取る.
  \item \textbf{定性的な結論:} 一般項を書かずに,収束,発散,周期,増減,
    誤差の大きさなどを示す.
\end{enumerate}

高校数学では第1段階が主役になりやすい.しかし,第2段階や第3段階が
妥協ではない場合も多い.大学数学では,むしろ第2段階の表現を選ぶこと
自体が問題の中心になる.この章のパターン集は,第1段階へ進む道だけで
なく,第2段階と第3段階へ進む道も見えるように構成する.
